\clearpage
\section{Разработка алгоритмов}
\subsection{Архитектура алгоритма построения графа знаний}

Разработанная система является конвейером для автоматического нахождения,
типизации и инкрементального сохранения семантических связей между
заметками (Markdown-файлами) в хранилище Obsidian.
Процесс построения состоит из нескольких шагов:
от сканирования хранилища до записи найденных отношений в файлы.
Алгоритм поддерживает сохранение состояния файлов в директории и
контрольные точки (checkpoints),
которые позволяют алгоритму продолжить работу с момента прерывания и
нацелены на экономию вычислительных ресурсов.
Алгоритм реализован на языке Python
(исходный код доступен в директории \texttt{src/kg\_builder}).
Общая схема конвейера представлена на рисунке~\ref{fig:pipeline}.

\begin{figure}[H]
    \centering
    \begin{tikzpicture}[
            node distance = 6mm and 30mm,
            block/.style = {
                    rectangle, rounded corners=3pt,
                    draw=black!60, fill=blue!8,
                    text width=72mm, minimum height=9mm,
                    align=center, font=\small
                },
            ioblock/.style = {
                    rectangle, rounded corners=3pt,
                    draw=black!40, fill=gray!10,
                    text width=72mm, minimum height=8mm,
                    align=center, font=\small\itshape
                },
            sideblock/.style = {
                    rectangle, rounded corners=3pt,
                    draw=black!50, fill=orange!12,
                    text width=36mm, minimum height=8mm,
                    align=center, font=\small
                },
            arrow/.style = {-Stealth, thick, black!70},
            dasharrow/.style = {Stealth-Stealth, dashed, thick, black!50},
            label/.style = {font=\scriptsize\itshape, text=black!60}
        ]

        % Main flow
        \node[ioblock]  (vault)    {Хранилище Markdown-файлов (vault)};
        \node[block,    below=of vault]  (scan)     {1. Обнаружение изменений\\{\scriptsize\normalfont VaultManager}};
        \node[block,    below=of scan]   (index)    {2. Разбивка на чанки и индексация\\{\scriptsize\normalfont VectorStore / LanceDB}};
        \node[block,    below=of index]  (retrieval){3. Гибридный поиск пар-кандидатов\\{\scriptsize\normalfont BM25 + векторный поиск}};
        \node[block,    below=of retrieval] (rerank){4. Переранжирование кандидатов\\{\scriptsize\normalfont Cross Encoder}};
        \node[block,    below=of rerank] (llm)      {5. Типизация связи\\{\scriptsize\normalfont LLM (локально или API)}};
        \node[block,    below=of llm]    (export)   {6. Запись связей в хранилище\\{\scriptsize\normalfont ExportService}};
        \node[ioblock,  below=of export] (result)   {Обновленные заметки с типизированными связями};

        % Side checkpoint block, aligned vertically between retrieval and llm
        \node[sideblock, right=of rerank] (meta)
        {Метаданные\\и чекпоинты\\{\scriptsize\normalfont MetadataManager}};

        % Main arrows
        \draw[arrow] (vault)     -- node[label, right]{markdown-файлы}    (scan);
        \draw[arrow] (scan)      -- node[label, right]{новые / обновленные файлы} (index);
        \draw[arrow] (index)     -- node[label, right]{чанки + эмбеддинги}  (retrieval);
        \draw[arrow] (retrieval) -- node[label, right]{пары-кандидаты}     (rerank);
        \draw[arrow] (rerank)    -- node[label, right]{отфильтрованные пары} (llm);
        \draw[arrow] (llm)       -- node[label, right]{тип связи + reasoning}(export);
        \draw[arrow] (export)    --                                           (result);

        % Dashed arrows to/from checkpoint block
        \draw[dasharrow] (meta.west) -- (rerank.east);
        \draw[dasharrow] ([yshift=4mm]meta.west) -- ([yshift=4mm]retrieval.east);
        \draw[dasharrow] ([yshift=-4mm]meta.west) -- ([yshift=-4mm]llm.east);

    \end{tikzpicture}
    \caption{Общая схема конвейера построения графа знаний}
    \label{fig:pipeline}
\end{figure}

\subsubsection*{Этап 1. Сканирование хранилища и обнаружение изменений (VaultManager)}
На этом этапе модуль \texttt{VaultManager} рекурсивно обходит директорию хранилища,
формируя список всех \texttt{markdown}-файлов.
Для каждого файла вычисляется SHA-256 хеш-сумма его содержимого для идентификации изменений.
Полученный хэш сравнивается с сохраненным с предыдущего запуска.
В результате файлы разбиваются на три категории: новые (отсутствующие в метаданных),
измененные (не совпадает хеш) и удаленные (есть в метаданных, но отсутствуют в директории).
Если метаданных нет, то все файлы считаются новыми.

Данный подход позволяет обрабатывать только измененные файлы,
а не все хранилище целиком.

\subsubsection*{Этап 2. Разбивка на чанки и индексация (VectorStore / LanceDB)}
Каждый новый или измененный файл проходит процесс индексации.
Эта процедура состоит из двух частей: обработки содержимого заметок и подготовки индексов.

Изначально текст каждого Markdown-файла может состоять из двух частей:
основного содержимого и блока ссылок, отделенного заголовком \texttt{link\_header},
если они были сгенерированы ранее.
Индексируется только основной текст заметки.

Текстовый индекс строится на основе лемматизированного содержимого файлов.
Лемматизация проводится при помощи модели \texttt{ru\_core\_news\_sm}
или \texttt{en\_core\_web\_sm} из \texttt{spaCy}~\cite{spacy_software}
для файлов с содержимым на русском или английском языке соответственно.
Из текстов извлекаются леммы всех алфавитных и числовых токенов.
Преобразованные тексты передаются в LanceDB~\cite{LanceDB_online}, где на них строится FTS-индекс,
использующий алгоритм BM25.
Данный подход позволяет увеличить точность поиска, нивелируя проблему, связанную с изменением форм слов,
в которых, например, могут быть изменены окончания или беглые гласные в корне,
что особенно распространено в русском языке.
Стоит отметить, что вместе с преобразованным текстом сохраняется и оригинальный.

Построение векторного индекса начинается с разбиения содержимого файла на перекрывающиеся
чанки. Используется рекурсивное посимвольное разбиение (\texttt{RecursiveCharacterTextSplitter} из LangChain).
Данная стратегия разбивает текст иерархически по заданным разделителям, расположенным в порядке их приоритета,
(символ новой строки, точка и т.д.), при этом не превышая заданный размер чанка
(по умолчанию 3,000 символов).
Эта стратегия сохраняет целостность предложений и структуру Markdown-файлов, что крайне важно
для вычисления эмбеддингов и формирования контекста для LLM.

Далее для каждого чанка следует вычисление плотного векторного
эмбеддинга при помощи модели из Sentence-Transformers~\cite{sentence_transformers_hf},
например популярной мультиязычной модели \texttt{BAAI/bge-m3}.

Полученные индексы объединяются друг с другом, а также с метаданными в виде пути к файлу и хеша,
после чего сохраняются в базу данных.
Хранение обоих форматов данных в одной таблице позволяет использовать гибридный поиск.
LanceDB позволяет регулировать баланс между векторным и полнотекстовым поиском при помощи параметра
\texttt{vector\_weight} (где $\texttt{vector\_weight} \in [0, 1]$): при \texttt{vector\_weight}~$= 1.0$ используется только ANN-поиск,
при \texttt{vector\_weight}~$= 0.0$~--- только BM25,
промежуточные значения используют оба поиска, сортируя результаты по унифицированной оценке релевантности.

\subsubsection*{Этап 3. Гибридный поиск пар-кандидатов (Retrieval)}

На данном этапе для каждого чанка выполняется поиск семантически близких чанков из других файлов.
Доступно три стратегии поиска:

\begin{itemize}
    \item \texttt{broad}, основанная на векторном поиске.
          Для каждого чанка текущего файла производится поиск $k$ ближайших фрагментов
          из других файлов.
    \item \texttt{strict}.
          Она основана на поиске прямых упоминаний сущностей в текстах заметок.
          Поиск производится по лемматизированному тексту.
    \item \texttt{combined} объединяет обе стратегии с дедупликацией. Результаты
          \texttt{broad} и \texttt{strict} проходят соответствующий им пайплайн извлечения и обработки,
          после чего фрагменты объединяются и дедуплицируются с помощью коэффициента Жаккара,
          вычисляемого по леммам.
          Результаты \texttt{strict} приоритизируются. То есть
          если фрагмент, извлеченный \texttt{strict}, будет иметь значение коэффициента с фрагментом,
          извлеченным \texttt{broad}, больше порога,
          то чанк от \texttt{broad} будет отброшен.
\end{itemize}


Особенности работы и недостатки каждой стратегии описаны в разделе 1.2.

Результатом поиска является \texttt{INITIAL\_RETRIEVAL\_K} фрагментов текста.

\subsubsection*{Этап 4. Переранжирование кандидатов (Cross-Encoder)}
Данный этап проходят только кандидаты, извлеченные векторным поиском,
поскольку среди них зачастую много ложноположительных.

Для снижения шума применяется дополнительная фильтрация кандидатов с помощью Cross Encoder.
Каждая пара (чанк источника, чанк цели) подается на вход Cross Encoder
(модель типа BAAI/bge-reranker-v2-m3~\cite{bge_reranker_v2})
с инструкцией-вопросом.
Cross Encoder вычисляет более точную оценку семантической связности.
Пары сортируются по убыванию оценки. В финальный список попадают \texttt{RERANKER\_TOP\_K} пар,
у каждой из которых оценка не ниже порогового значения.

Все прошедшие фильтрацию чанки одной заметки группируются по файлу-цели.
На их основе создаются структуры \texttt{CandidatePair},
содержащие исходный и целевые чанки, а также пути к соответствующим файлам и оценку близости
из LanceDB и Cross Encoder.

Результатом этапа является список структур \texttt{CandidatePair},
к которым добавлены результаты лексического поиска, если он был использован.

\subsubsection*{Этап 5. Типизация связи (LLM)}
Прошедшие фильтрацию сгруппированные пары-кандидаты
передаются в языковую модель для определения точного типа семантической связи.
Для каждой такой пары отбираются не более \texttt{LLM\_TOP\_K\_CTX}
наиболее релевантных чанков для формирования контекста
на основе оценки Cross Encoder.
Важно уточнить, что результаты стратегии \texttt{strict}
являются более приоритетными, а
результаты стратегии \texttt{broad}
служат для дополнения выборки и
добавляются в итоговый список
по остаточному принципу на основе их оценок.

Логика запросов к LLM поддерживает параллелизм и батчинг.
Количество одновременных запросов к модели ограничено заданным числом и контролируется семафором.
Кроме того, в рамках одного запроса может проводиться типизация сразу нескольких ссылок.
Это значительно снижает накладные расходы и время обработки.

Важно отметить обработку ошибок: у каждого батча есть три попытки, при израсходовании которых
батч откладывается и логируется ошибка. Он будет повторно обработан при следующем запуске.

В рамках запроса модель получает системный промпт со списком допустимых типов связей и инструкциями,
а также пользовательский промпт с парами файлов и их контекстом.

Ответ модели возвращается в формате JSON. Его парсинг производится посредством JsonOutputParser
из LangChain, позволяющий восстанавливать синтаксически некорректный JSON.
При возврате моделью раздела с размышлениями
он логируется и вырезается для упрощения парсинга.

Для работы с LLM используется LangChain,
который позволяет интегрировать
как локальные модели
(например, запущенные через llama.cpp~\cite{Gerganov_llamacpp}),
так и облачные провайдеры (OpenAI, Google, и другие).

После каждого батча предсказания модели сохраняются в файл,
формируя контрольную точку, что позволяет возобновить типизацию
в случае ошибки или аварийного завершения программы.
Данная оптимизация призвана
убрать необходимость повторять выполненную работу.

\subsubsection*{Этап 6. Сохранение связей (ExportService)}
На заключительном этапе типизированные связи сохраняются в одном из трех режимов:

\begin{itemize}
    \item \texttt{inplace}.
          Связи добавляются непосредственно в исходные Markdown-файлы под специальным заголовком
          \texttt{LINK\_HEADER}.
          При наличии связей с предыдущих запусков
          новые связи добавятся к уже имеющимся, не создавая дубликатов.
    \item \texttt{json}. Результаты экспортируются в JSON-файл без изменения исходных заметок.
          Данный режим создан для удобного переиспользования результатов, например, для вычисления метрик.
    \item \texttt{export}. Создается полная копия хранилища с вписанными связями,
          исходные файлы остаются нетронутыми.
          Этот формат сохранения полезен для экспериментов,
          таких как оценка GraphRAG, требующих полноценного хранилища для работы.
\end{itemize}

Формат ссылки конфигурируется
(по умолчанию используется синтаксис Dataview).

\subsubsection*{Управление состоянием и отказоустойчивость (MetadataManager)}

При первом запуске конвейера стандартные параметры из \texttt{config.py}
сохраняются в \texttt{config.json} внутри директории \texttt{.kg\_builder/}.
Последующие запуски считывают этот файл и выставляют настройки, заданные в нем,
что обеспечивает воспроизводимость и
повышает удобство использования на разных хранилищах с индивидуальными параметрами.

Опция \texttt{--ignore-local-config} заставляет пайплайн
проигнорировать \texttt{config.json} и использовать значения из \texttt{config.py}.

Алгоритм сохраняет свое состояние с помощью создания снимка, который включает
текущий этап,
его прогресс и прочие метаданные, необходимые для возобновления работы.

Наибольшей детализацией сохранения обладают этапы ранжирования кандидатов и LLM-классификации отношений.
У них есть логика контрольных точек, при которых сохраняется список уже обработанных файлов или пар-кандидатов.
При возобновлении работы пайплайна на одном из этих шагов в алгоритм будут переданы только необработанные данные.

При последующих запусках \texttt{KnowledgeGraphBuilder} проверяет последнее сохраненное состояние.
Если предыдущий запуск был прерван, то он продолжит обработку с того же шага, минуя уже выполненные.

Это сохраняет значительное количество вычислительных ресурсов и времени, особенно на дорогостоящем LLM-этапе.

Все метаданные сохраняются в отдельной директории \texttt{.kg\_builder/.out/}:
\begin{itemize}
    \item \texttt{candidates.json}~--- начальные кандидаты до переранжирования;
    \item \texttt{reranked\_candidates.json}~--- кандидаты после переранжирования,
          обогащенные результатами LLM (поля \texttt{llm\_type}, \texttt{llm\_reasoning}
          дописываются после классификации)
    \item \texttt{predictions\_partial.json}~--- предсказания LLM, обновляемые инкрементально
    \item \texttt{run\_state.json}~--- текущее состояние алгоритма
\end{itemize}

Специальный флаг \texttt{--fresh-start} позволяет перезаписать метаданные,
в том числе конфиг и список обработанных файлов.
